\chapter*{Tóm tắt}
\label{summary}

Sự phát triển mạnh mẽ của các mô hình ngôn ngữ lớn (LLM) đã mở ra tiềm năng ứng dụng trí tuệ nhân tạo vào lĩnh vực giảng dạy lập trình. Tuy nhiên, các công cụ AI hiện tại như ChatGPT hay GitHub Copilot thường cung cấp đáp án trực tiếp, vi phạm nguyên tắc sư phạm và tạo ra sự phụ thuộc của người học vào AI thay vì phát triển tư duy giải quyết vấn đề. Bên cạnh đó, các LLM thương mại chưa được tinh chỉnh cho nhiệm vụ gỡ lỗi code vẫn còn nhiều hạn chế về độ chính xác, đặc biệt ở khả năng sửa lỗi tự động (Code Repair). Đề tài này hướng đến xây dựng hệ thống \textit{Socrates Code Tutor} --- một hệ sinh thái gia sư lập trình AI tích hợp phương pháp Socratic, bao gồm VS Code Extension, Backend API, Web Portal và hệ thống chấm bài tự động. Về mặt kỹ thuật, đề tài áp dụng kỹ thuật tinh chỉnh tham số hiệu quả LoRA (Low-Rank Adaptation) để huấn luyện ba mô hình ngôn ngữ lớn --- Gemma-4-E4B-IT, Qwen2.5-Coder-3B-Instruct và Qwen2.5-3B-Instruct --- trên hai bộ dữ liệu là DebugEval (24.892 mẫu) và Socratic Debugging Dataset (552 mẫu huấn luyện); AI Runtime được xây dựng theo kiến trúc LangGraph với các workflow chuyên biệt và cơ chế chống tiết lộ đáp án đa tầng (Anti-Solution Guard). Kết quả thực nghiệm cho thấy mô hình Qwen2.5-Coder sau tinh chỉnh đạt mức cải thiện đồng đều trên cả bốn nhiệm vụ DebugEval (Bug Localization tăng 27 điểm phần trăm, Bug Identification tăng 18 điểm phần trăm), trong khi mô hình Socratic Tutor đạt tỷ lệ đặt câu hỏi dẫn dắt 96,1\% và tỷ lệ che giấu đáp án tuyệt đối 100\%, vượt mục tiêu đề ra. Hệ thống đề xuất không chỉ đóng góp bằng chứng thực nghiệm về hiệu quả của LoRA trong lĩnh vực giáo dục lập trình mà còn cung cấp một nền tảng mã nguồn mở có thể triển khai thực tế tại các cơ sở giảng dạy đại học.

\bigskip
\noindent\textbf{Từ khóa:} gia sư lập trình AI, phương pháp Socratic, tinh chỉnh mô hình ngôn ngữ lớn, LoRA, gỡ lỗi code, VS Code Extension, LangGraph.
